\documentclass[conference]{IEEEtran}
\usepackage[T1]{fontenc}
\usepackage{cite}
\usepackage{url}
\usepackage{array}
\begin{document}
\title{From Scores to Supported Claims: Measurement Audits of Document and Retrieval Evaluation}
\author{}
\maketitle
\begin{abstract}
Evaluation scores are used to compare systems, but their interpretation
depends on what an evaluator observes, how observations are scored and
what comparison is supported. We examine these dependencies in three
audits: a staged document evaluator, the released FinRank
evidence-ranking benchmark and researcher-generated runs on two BEIR
datasets. In the staged evaluator, 96 of 118 controlled header changes
were visible; under a single-responsive-leaf restriction, flat and
hierarchical scores had the same ordering up to an approximately 0.375
scale factor, despite different fixed-threshold detection rates. In
FinRank, changing the tie policy shifted absolute pairwise accuracy
while preserving the direction of the examined hard-versus-random
comparison. For the BEIR runs, a paired sign-flip analysis with Holm
correction supported two of ten lexical-system comparisons on SciFact
and one on NFCorpus. An exact-tie-aware point-estimate baseline also
passed our duplicate and degradation controls. These findings show how
explicit measurement contracts qualify or preserve particular
comparisons; they do not establish a universally validated audit method.
\end{abstract}
\begin{IEEEkeywords}
measurement validity, document evaluation, information retrieval,
reproducibility, benchmark auditing.
\end{IEEEkeywords}
\section{Introduction}\label{introduction}

A lower detection rate can indicate a weaker evaluator, but it can also
arise when the same numerical threshold is applied to scores on
different scales. Likewise, a change in benchmark score need not change
the comparison for which the score was reported; and the highest mean in
a leaderboard does not by itself establish that every apparent ordering
is supported by the data. Such distinctions matter when researchers use
evaluation results to select systems or interpret a benchmark.

We ask: \textbf{which interpretations remain supported when the
observation surface, scoring rule and comparison procedure are made
explicit?} We investigate three evaluation constructions. First,
controlled changes to public financial tables probe a staged document
evaluator. Second, we inspect tie handling and source identity in
FinRank \cite{r1}. Third, we assess researcher-generated lexical
retrieval runs on SciFact and NFCorpus from BEIR \cite{r2}, with a later
dense-model exploration. The BEIR runs are our own constructions, not
the published benchmark leaderboard.

Our contributions are empirical and bounded. We explain a detection-rate
difference by a positive score rescaling under a single-responsive-leaf
restriction; show that a particular FinRank comparison retains its
direction despite changes in absolute scores; and distinguish
point-estimate rankings from comparisons supported under a specified
paired statistical procedure. Each finding is tied to an explicit
measurement object and relevant limitations.

Prior work has examined metric and construct validity \cite{r3},
\cite{r4}, \cite{r5}. AgentSuite \cite{r6} audits interacting components
in agent benchmarks at broader scale. Our study instead examines
arithmetic, observation and comparison contracts in specific document
and retrieval evaluations. It does not propose a new general theory of
validity or claim an observed change in a user's decision.

Document-specific work also evaluates semantic alignment and stage-level
extraction: SCORE \cite{r7} addresses valid alternative representations
in generative parsing, while Soric et al.~\cite{r8} benchmark table
extraction across heterogeneous PDFs and analyze evaluation metrics. Our
case studies ask a complementary question about the interpretation of an
existing evaluation's thresholds and retrieval comparisons; we do not
benchmark document extraction models or claim a superior parsing metric.

\section{Audit Design}\label{audit-design}

We distinguish three questions: whether a change is \textbf{visible} at
a selected evidence surface; whether a specified predicate
\textbf{detects} it; and whether it changes an outcome under a stated
equivalence contract. Visibility, detection and outcome change are
separate observations. We then examine score sensitivity, the stability
of a named comparison and numerical reproducibility.

The designs differ. Case 1 uses controlled perturbations. Case 2
inspects a released benchmark, including source-informed post-hoc
identity analysis. Case 3 uses researcher-generated systems and test
queries from a public benchmark; its original audit choices were not
independently timestamped before the first run. We preserve fixed
outputs and identify later explorations rather than treating them as
prospective confirmation.

The cases target three interpretations. First, a detection-rate
difference does not necessarily imply a ranking difference when the
score scales differ. Second, a shift in absolute accuracy need not
reverse a stated comparison. Third, a higher retrieval mean is a point
estimate; it does not automatically establish every pairwise system
ordering. The following table summarizes what the audit changes in each
interpretation; the cases below specify what is held fixed and what is
varied.

\begin{table}[htbp]
\caption{Interpretations qualified or retained by the three audits.}
\label{tab:interpretations}
\centering
\footnotesize
\renewcommand{\arraystretch}{1.12}
\begin{tabular}{@{}>{\raggedright\arraybackslash}p{.18\columnwidth}>{\raggedright\arraybackslash}p{.34\columnwidth}>{\raggedright\arraybackslash}p{.36\columnwidth}@{}}
\hline
Case & Tempting interpretation & Supported interpretation \\
\hline
Document & Lower detection means weaker ranking & Rates differ; ranks coincide in the single-leaf setting. \\
FinRank & Absolute score changes imply a changed comparison & The random-minus-hard gap stays positive under both inspected tie policies. \\
Retrieval & A higher mean establishes every system ordering & The paired test supports some differences; other pairs are not proven equal. \\
\hline
\end{tabular}
\end{table}

For each construction we identify the input and outcome carrier, the
evaluator or scoring implementation, the original reporting convention,
the comparison of interest, and any alternative rule applied to fixed
outputs. We use the term \emph{decision} for a stated analytical
interpretation such as the direction of a comparison; no intervention
with a human decision-maker was conducted. The analysis can preserve an
initial interpretation as well as qualify it. An audit of a score is
useful even when the direction remains stable, because it identifies
which aspects of the scoring contract matter to that conclusion.

The cases are complementary rather than replications of one universal
intervention. The document case compares response visibility and
threshold behavior in an instrumented construction. The FinRank case
varies tie credit for a released financial retrieval benchmark. The BEIR
case investigates metric definitions, query-level variation and
uncertainty for newly constructed runs. Consequently, their sample
counts, populations and effect sizes are not combined into one overall
audit-success rate.

\section{Controlled Document
Evaluation}\label{controlled-document-evaluation}

The frozen staged evaluator has parsing and question-answering
checkpoint trees with twelve scored leaves. Its seed-42 panel contains
95 TAT-QA-derived source tables \cite{r9} and 118 controlled header
changes across six families. Mapping and regression artifacts form the
observed surface. For each injected header, the visibility analysis
identifies its original canonical field: the change is counted as
absorbed when that field still maps to the injected header, and visible
otherwise. This matcher-relative criterion is not simply a
character-level difference. Protocol detection and downstream harm are
separate quantities. The audit holds the evaluated inputs and scores
fixed while comparing declared aggregation and threshold conventions. It
does not observe every internal state or possible downstream use.

Of 118 injections, 96 were visible. None of the 13 OCR-like changes was
visible on this surface. This observation does not establish 13 harmful
misses: one existing semantic judge classified the cases as
referentially equivalent, without an independent human reference.

In this construction the QA tree was frozen and only the field-hit-rate
leaf responded. The unrounded hierarchical contrastive score was a
positive rescaling of the flat score, with fitted slope approximately
0.375 and identical observed rankings. The original flat field-hit-rate
drop and hierarchical parsing-root drop each used a numerical operating
threshold of \textbf{at least 2.0}, on their respective scales. Under
the observed scaling, a flat threshold of 2.0 corresponds to about 0.75
on the hierarchical scale, rather than 2.0. At these original operating
points the flat and hierarchical protocols detected 75/95 versus 20/95
tables, or 85/118 versus 13/118 injections. Tables and injections are
distinct units. The difference in rates therefore does not by itself
establish a difference in ranking ability.

\begin{table}[t]
\caption{Original seed-42 operating-point detection; table and injection counts are different units.}
\label{tab:original-detection}
\centering
\begin{tabular}{lrr}
\hline
Unit & Flat & Hierarchical \\
\hline
Source tables & 75/95 & 20/95 \\
Header injections & 85/118 & 13/118 \\
\hline
\end{tabular}
\end{table}

The separate matched-threshold reanalysis used strict greater-than
comparisons: a flat threshold of 1.0 corresponds to a hierarchical
threshold of 0.375; the recorded fold detects 21 cases under either. It
is not the original 2.0 operating-point analysis. This distinction also
matters for boundary values: 12 of the 95 tables have a flat drop
exactly equal to 2.0.

This comparison illustrates why a threshold needs both a numeric value
and a scale. In the observed single-leaf setting, multiplying every flat
contrast by a positive constant preserves order. A corresponding
threshold must be multiplied by the same constant if the same decision
rule is to be retained. The original protocols did not apply
corresponding thresholds. The matched-grid example tests the
relationship on saved outputs; it does not certify that a production
threshold has been calibrated. Neither the scaling result nor the
matched-grid behavior should be generalized to a setting where
additional leaves respond differently.

The clean-only contrastive construction compares repeated deterministic
computations on an unchanged record. All 2,382 usable clean records have
zero scores, so every positive threshold has observed clean
false-positive rate zero; this panel does not identify a unique
deployment threshold. Across three seeds, the flat descriptive table-run
ratio was 219/282 using 119 unique source-table identifiers. A post-hoc
UID-cluster bootstrap gives a 95\% interval of {[}0.7117, 0.8364{]}; the
repeated table runs are not 282 independent source tables.

The per-seed contributions to the 219/282 ratio are 75/95, 76/97 and
68/90. The denominators refer to table appearances within each run; the
same source table may contribute to more than one run. The bootstrap
samples source identifiers with replacement, includes all run
appearances for each sampled identifier, and recalculates the pooled
ratio. This uncertainty interval is post-hoc and conditional on the
available source tables and three seeds; it should not be interpreted as
a guarantee for future datasets or seeds. Equal weighting of source
identifiers gives 0.7535 with interval {[}0.6835, 0.8207{]}.

In a separate stage-aligned control, 72 source clusters contributed one
injection each for mapping, retrieval and numeric processing. A
leave-one-cluster-out nearest-centroid classifier on the full response
vector attained macro-F1 and top-1 1.000. The three directly targeted
leaves gave the same result; removing them reduced top-1 to 0.333 under
the frozen tie-as-incorrect rule, with a 0.667 tie rate. The observed
stage separation is attributable to explicit instrumentation, not to
hierarchy alone. A nested FinQA probe \cite{r10} covered 89 rows from 30
report clusters; it changed program fields but did not recompute
downstream result fields. Its zero detected and zero changed outcomes
cannot estimate a failure rate for changed answers.

\section{FinRank: Score Sensitivity and Decision
Stability}\label{finrank-score-sensitivity-and-decision-stability}

FinRank ranks evidence from financial filings \cite{r1}. Under
tie-as-positive scoring, a score tie counts for the gold passage; under
strict-win scoring, it does not. The tie comparisons and materiality
threshold were fixed before pair-level outcomes. The following table
shows that random-negative accuracy moves by 1.084 percentage points for
TF-IDF and 1.029 for BM25, exceeding the fixed 1.0-point threshold,
while the random-minus-hard gap remains positive in both systems.

\begin{table}[htbp]
\caption{FinRank pairwise accuracy under two tie policies; gaps are computed before rounding.}
\label{tab:finrank}
\centering
\begin{tabular}{llrrr}
\hline
System & Policy & Hard (\%) & Random (\%) & Gap (pp) \\
\hline
TF-IDF & Positive & 77.756 & 91.935 & 14.179 \\
TF-IDF & Strict & 77.030 & 90.850 & 13.821 \\
BM25 & Positive & 74.780 & 87.789 & 13.009 \\
BM25 & Strict & 74.094 & 86.760 & 12.666 \\
\hline
\end{tabular}
\end{table}

Gaps are random accuracy minus hard accuracy, computed before display
rounding. The comparison's qualitative direction is stable under these
two policies; this does not establish stability for all decisions. A
later source-informed identity analysis reproduces 442 hard-negative
occurrences across 305 records sharing canonical text with another
record's gold passage, and gold-passage loss in 33 records under the
inspected metadata filter. These are post-hoc reproductions of known
properties, not newly discovered benchmark defects. In-record and
pairwise metrics reproduce exactly in the audited environment; global
and metadata-filtered metrics differ from the released results by up to
0.19 and 0.16 percentage points, respectively, with the cause
unresolved.

The identity analysis also explains why text identity and record
identity are different contracts. A pooled passage can occur in several
records with different metadata. When the inspected implementation
deduplicates text by first occurrence and retains that occurrence's
metadata, a later record's gold text can be excluded by metadata
filtering even though its text appeared in the pooled corpus. The
reproduced counts describe that implementation and its released data,
not the prevalence of this behavior in other retrieval systems. Exact
in-record reproduction and the remaining global mismatch are reported
separately because one does not resolve the other.

\section{Retrieval Comparisons}\label{retrieval-comparisons}

We built five lexical configurations on BEIR SciFact and NFCorpus
\cite{r2}: TF-IDF and BM25 parameter pairs (k1, b) of (0.9, 0.4), (1.2,
0.75), (0.5, 0.3) and (1.5, 0.6). We concatenate each document's title
and text, lowercase and tokenize on ASCII alphanumeric runs, use
log-scaled term frequency and smoothed inverse document frequency for
TF-IDF cosine, and resolve equal scores by ascending document ID unless
testing the alternative order. They use the released corpora, qrels and
300/323 judged queries, respectively. The systems are our own runs, and
the two datasets belong to one benchmark family. We compare system
scores on the same queries so that statistical analysis is paired; each
dataset has its own comparison family.

Our initial scoring implementation binarized graded relevance and used a
different MAP denominator. Aligning it with the trec\_eval contract
corrected both definitions. Under graded nDCG@10, TF-IDF moved from
second to third in one NFCorpus ranking; the separate denominator
correction changed TF-IDF's MAP@10 from approximately 0.2200 to 0.1162.
This is a correction of our implementation, not evidence that the
official leaderboard was wrong. For fixed rankings, trec\_eval 10.0-rc3
agreed with 18,690 per-query nDCG, MAP and recall values within 1e-4; it
does not validate tie-breaking policy.

The cross-check writes a strictly decreasing score for each fixed rank,
so the reference evaluator receives the order selected by our run. It
can test the metric values for that order, but cannot decide which
document should win an equal-score tie. The maximum observed discrepancy
was approximately 0.00005, consistent with four-decimal output
precision. The executable used was trec\_eval itself; the check was not
a separate execution of the pytrec\_eval Python binding. These
distinctions make the reproducibility statement narrower and more useful
than a general assertion that an entire retrieval pipeline has been
independently reproduced.

For the five systems, 2,000 paired bootstrap resamples produced
intervals excluding zero for 3/10 pairs on SciFact and 4/10 on NFCorpus.
A separate two-sided paired sign-flip procedure with 2,000 Monte Carlo
draws per pair and Holm correction supported 2/10 and 1/10,
respectively. These are different inference procedures, so the count
difference is not attributable to Holm correction alone. The inference
depends on the relevant exchangeability or symmetry and
query-independence assumptions; non-significance does not establish
equality.

The supported SciFact pairs were BM25(1.5, 0.6) versus TF-IDF (mean
nDCG@10 difference +0.03334; Holm-adjusted p = 0.0100) and BM25(1.2,
0.75) versus TF-IDF (+0.02930; p = 0.0450). On NFCorpus, BM25(1.5, 0.6)
versus BM25(0.5, 0.3) had a difference of +0.00817 (p = 0.0300). A later
200,000-draw sensitivity check on regenerated lexical query scores
retains the same three supported pairs when Holm is recalculated within
the five-system family. It reduces Monte Carlo variability but does not
independently validate the data or inferential assumptions.

Separately from varying cutoff k, reversing document-ID tie order for
one NFCorpus BM25 configuration at k=10 changed the ordered top ten for
36.5\% of queries and top-ten membership for 25.7\%. Mean nDCG@10
changed from 0.30453945 to 0.30449093, approximately 0.00005.
Query-level changes can therefore be hidden by the aggregate in this
setting; this is not an estimated benchmark-wide failure rate.

For an intentionally duplicated system and an intentionally degraded
control, an exact-tie-aware point-estimate rule and the paired audit
both passed both controls. The controls do not demonstrate that the
audit outperforms the simpler rule. A subsequent BGE-small-en-v1.5
\cite{r11} exploration gave mean nDCG@10 of 0.7127 versus 0.6622 for
BM25 on SciFact, and 0.3438 versus 0.3062 on NFCorpus. This is an
exploratory retrieval result selected after the original analysis, not
evidence that the audit improved a user decision or that lexical tie
findings transfer to dense models.

The BGE comparison was later included with the five lexical runs in a
separate six-system family of 15 pairs per dataset. With 200,000 Monte
Carlo sign flips per pair and Holm correction within each dataset, the
BGE-minus-BM25 nDCG@10 difference was +0.05049 on SciFact (adjusted p =
0.066000) and +0.03759 on NFCorpus (adjusted p = 0.015540). These
post-hoc tests are not the original ten-pair family or independent
confirmation. The saved dense outputs contain ranked document IDs rather
than original dense scores, so they cannot establish dense-score tie
frequencies. Numerical reproducibility of the tests does not establish
their inferential assumptions.

The single BGE comparison originally used a paired bootstrap: the mean
differences were 0.0505 {[}0.0137, 0.0891{]} on SciFact and 0.0376
{[}0.0154, 0.0608{]} on NFCorpus, uncorrected for a comparison family.
It used normalized embeddings, the model's retrieval query prefix and a
512-token cap. A separate trec\_eval execution agreed on 1,246 dense and
BM25 per-query scores to the reference program's four-decimal print
precision. The later 15-pair test answers a different question about
multiplicity and must not be represented as a prospective verification
of the selected BGE comparison.

\section{Discussion and Conclusion}\label{discussion-and-conclusion}

The three cases show distinct ways in which an evaluation contract
constrains a claim. In the staged evaluator, a rate difference coexists
with identical ranks under a restricted score rescaling. In FinRank, a
tested comparison retains direction as absolute scores change. In our
BEIR runs, metric definitions, ties and statistical comparison
procedures affect the interpretation of rankings.

These findings have limited scope: the first case uses controlled table
perturbations, FinRank uses lexical baselines and has an unresolved
reproduction gap, and the BEIR analyses are retrospective runs over two
datasets with a post-hoc dense extension. Existing validity frameworks
already cover broader audit principles \cite{r3}, \cite{r4}, \cite{r5}.
Our contribution is a reproducible set of case-specific observations and
the distinction between numerical movement and supported interpretation.
We recommend reporting the measurement object, scoring convention,
comparison and uncertainty together; we do not claim a general audit
guarantee or observed deployment benefit.

\section{Data and Materials}\label{data-and-materials}

The TAT-QA, FinRank and BEIR data used in this study are described in
\cite{r9}, \cite{r1} and \cite{r2}, respectively; the BEIR runs use
released relevance judgments. Study-specific audit scripts,
configuration records and aggregate results are retained as research
artifacts. Access to and redistribution of underlying data remain
subject to the respective source terms.

\section*{Acknowledgment}

Generative AI tools, including DeepSeek \cite{r12} and OpenAI Codex
\cite{r13}, assisted with analysis-code development and drafting and
revision of manuscript text in Sections I--VI. Google Gemini \cite{r14}
provided review suggestions on earlier drafts. The authors are
responsible for the final manuscript and its conclusions.
\begin{thebibliography}{14}
\bibitem{r1} S. Mansouri et al., ``FinRank: An Evidence-Grounded Benchmark for
Financial Question Answering and Retrieval over SEC Filings,''
arXiv:2608.07400, 2026.
\bibitem{r2} N. Thakur et al., ``BEIR: A Heterogeneous Benchmark for Zero-shot
Evaluation of Information Retrieval Models,'' NeurIPS Datasets and
Benchmarks, 2021.
\bibitem{r3} C. Ross et al., ``What makes a good metric? Evaluating automatic metrics
for text-to-image consistency,'' arXiv:2412.13989, 2024.
\bibitem{r4} W. Caban, ``Measurement Without Validity: The Compounding Reliability
Problem in Agentic AI Evaluation,'' arXiv:2608.00794, 2026.
\bibitem{r5} A. M. Bean et al., ``Measuring what Matters: Construct Validity in Large
Language Model Benchmarks,'' arXiv:2511.04703, 2025.
\bibitem{r6} H. Suh et al., ``AgentSuite: Toward More Reliable Agent Evaluation with
a Component-Based Benchmark Auditing Pipeline,'' ICML, PMLR
306:116300--116319, 2026.
\bibitem{r7} R. Li et al., ``SCORE: A Semantic Evaluation Framework for Generative
Document Parsing,'' arXiv:2509.19345, 2025.
\bibitem{r8} M. Soric, C. Gracianne, I. Manolescu, and P. Senellart, ``Benchmarking
Table Extraction from Heterogeneous Scientific PDF Documents,''
arXiv:2511.16134v2, 2026.
\bibitem{r9} F. Zhu et al., ``TAT-QA: A Question Answering Benchmark on a Hybrid of
Tabular and Textual Content in Finance,'' ACL-IJCNLP, 2021, doi:
10.18653/v1/2021.acl-long.254.
\bibitem{r10} Z. Chen et al., ``FinQA: A Dataset of Numerical Reasoning over Financial
Data,'' EMNLP, 2021, doi: 10.18653/v1/2021.emnlp-main.300.
\bibitem{r11} Beijing Academy of Artificial Intelligence, ``bge-small-en-v1.5,'' model
card, revision 5c38ec7c, https://huggingface.co/BAAI/bge-small-en-v1.5.
\bibitem{r12} DeepSeek, ``DeepSeek V4.1,'' https://www.deepseek.com/, accessed
November 2, 2026.
\bibitem{r13} OpenAI, ``Codex,'' https://openai.com/codex/, accessed November 2, 2026.
\bibitem{r14} Google, ``Gemini,'' https://gemini.google.com/, accessed November 2,
2026.
\end{thebibliography}
\end{document}
